\documentclass[12pt]{article}

\usepackage{fontawesome}
\usepackage{hyperref}
\usepackage{xurl}
\usepackage[tagged, highstructure]{accessibility}
\usepackage{bookmark}
\usepackage{enumitem}

\hypersetup{
    colorlinks=false,
    pdfborder={0 0 0},
    pdftitle={Networks and permissions},
    pdfauthor={Adrianna Holden-Gouveia},
    pdfsubject={Introduction to Linux - Lab Assignment},
    pdflang={en-US},
    pdfkeywords={lab assignment, Networks and permissions}
}

% Configure PDF bookmarks for navigation
\bookmarksetup{
    numbered,
    open,
}

% Configure list spacing for better accessibility
\setlist{nosep}



\title{Networks and permissions}
\author{
        Adrianna Holden-Gouveia \\
        Website: \href{https://aholdengouveia.name}{aholdengouveia.name}\\ 
        \date{\vspace{-5ex}}
        %Email: \href{mailto:admin@aholdengouveia.name}{admin@aholdengouveia.name} \\
%        \faLinkedin{: aholdengouveia} \\
        \faGithub {: aholdengouveia} \\
%        \faTwitter {: aholdengouveia} \\
        }

%S\date{\today}


\begin{document}    

\maketitle

\section*{Accessibility Notice}
This document is also available in HTML format at:

Link: \href{https://aholdengouveia.name/IntroLinux/labs/networking.html}{\textbf{https://aholdengouveia.name/IntroLinux/labs/networking.html}}

The HTML version provides enhanced accessibility features including keyboard navigation, screen reader support, responsive design, dark mode support, and high contrast options.


%\begin{abstract}
%This is a template for Linux Administration Lab work
%\end{abstract}
%\tableofcontents

\section*{Objectives:}
\begin{itemize}
    \item Learn basic networking commands (ping, ifconfig/ip)
    \item Understand how to test network connectivity
    \item Master Linux file permissions and the chmod command
    \item Practice setting different permission levels for files
\end{itemize}
\section*{Complete the following problems}

References, a video, a PowerPoint and some notes are available at my website
Link: \href{https://www.aholdengouveia.name/IntroLinux/networking.html}{\textbf{https://www.aholdengouveia.name/IntroLinux/networking.html}}

\subsection*{Networking}
You are going to use two basic commands ping and ifconfig or ip a

    \begin{enumerate}
        \item ping Google.com Was ping successful? What was the result? Take a screenshot
        \item ping CNN.com Was ping successful? What was the result? Take a screenshot
        \item Use ifconfig or ip a What does it show you? Use man to see what other useful things one of those can do
    \end{enumerate}

\subsection*{File Permissions}

Create a subdirectory for this lab and put several files in it either using cp or touch


Use chmod to create files with the following permissions:
    \begin{enumerate}
        \item Everyone can read and write but not execute
        \item Owner has full permissions but everyone else can only read
        \item No one can do anything including the owner
        \item Owner and group have read and write permission, public only has read
        \item Owner can read and execute but not write, everyone else can read\\
    \end{enumerate}


Use ls -l to show results, take a screenshot and include in report

Pick one file and explain the permissions in your own words in the lab report


\section*{Deliverables}
 One file with your name, date, and lab title at the top and all the answers. 
 \begin{itemize}
        \item You should have screenshots for each ping, your ifconfig or ip a, and showing your file permissions
        \item Make sure to include labels for each screenshot so it's clear what you're turning in. \item Don't forget to include the answers to each question asked as well. 
\end{itemize}
\end{document}